\documentclass[11pt]{article}
\usepackage[utf8]{inputenc}
\usepackage{amsmath,amssymb,booktabs,geometry,verbatim}
\usepackage{CJKutf8}
\geometry{margin=1in}

\title{CSP-J 2024 (Junior)}
\date{2024-10-26, 8:30--12:00 (3.5 hours)}

\begin{document}
\maketitle

\section*{Problem 1: Poker Cards (\begin{CJK}{UTF8}{gbsn}扑克牌\end{CJK})}

Time limit: 1000 ms.
Memory limit: 512 MB.

Little P borrowed a deck of $n$ poker cards from his classmate Little Q.

In this problem we do not consider the jokers.
Then every card has two attributes: a suit and a rank.
There are $4$ suits: diamonds, clubs, hearts and spades.
There are $13$ ranks; from smallest to largest they are
\texttt{A 2 3 4 5 6 7 8 9 T J Q K}.
Note: in this problem the rank $10$ is written as \texttt{T}.

We call a deck of poker cards \textbf{complete} if and only if for every
suit and every rank there is \textbf{exactly} one card with that suit and
that rank.
Hence a complete deck has exactly $4 \times 13 = 52$ cards.
The following picture shows all 52 cards of a complete deck.

(The picture, omitted here, displays the 52 cards arranged in four rows,
one row per suit, each row running through the 13 ranks.)

The deck Little P borrowed may not be complete, so Little P plans to borrow
some more cards from his classmate Little S.
You may assume that Little S has infinitely many cards of every kind, so
Little P can choose the borrowed cards arbitrarily.
Little P wants to know the minimum number of cards he must borrow from
Little S so that, from the cards borrowed from Little S and from Little Q,
one can select $52$ cards that form a complete deck.

To make the input convenient, we use the character \texttt{D} for diamonds,
\texttt{C} for clubs, \texttt{H} for hearts and \texttt{S} for spades.
In this way every card can be represented by a string of length $2$, whose
first character gives the suit of the card and whose second character gives
its rank.
For example, \texttt{CA} denotes the club \texttt{A} and \texttt{ST}
denotes the spade \texttt{T} (the spade 10).

\subsection*{Input}

The first line of input contains an integer $n$, the number of cards.

The next $n$ lines follow:

Each line contains a string of length $2$ describing one card, whose first
character describes its suit and whose second character describes its rank.

\subsection*{Output}

Output one line containing one integer: the minimum number of cards that
still have to be borrowed from Little S in order to assemble a complete deck.

\subsection*{Sample 1 input}

\begin{verbatim}
1
SA
\end{verbatim}

\subsection*{Sample 1 output}

\begin{verbatim}
51
\end{verbatim}

\subsection*{Explanation of sample 1}

This deck contains one spade \texttt{A}.
Little P still has to borrow the 51 cards other than the spade \texttt{A}
to form a complete deck.

\subsection*{Sample 2 input}

\begin{verbatim}
4
DQ
H3
DQ
DT
\end{verbatim}

\subsection*{Sample 2 output}

\begin{verbatim}
49
\end{verbatim}

\subsection*{Explanation of sample 2}

This deck contains two diamond \texttt{Q}, one diamond \texttt{T}
(the diamond 10) and one heart 3.
Little P still has to borrow the $49$ cards other than the heart 3, the
diamond \texttt{T} and the diamond \texttt{Q}.

\subsection*{Sample 3}

See \texttt{poker/poker3.in} and \texttt{poker/poker3.ans} in the
contestant's directory.
(On Luogu these files are provided in the attachment \texttt{poker.zip}.)

\subsection*{Explanation of sample 3}

This deck is complete, so no further card needs to be borrowed.

This sample satisfies the condition that all cards are given in increasing
order of rank, and cards of equal rank are given in the order diamonds,
clubs, hearts, spades.

\subsection*{Constraints}

For all test data it is guaranteed that $1 \leq n \leq 52$, and that each of
the $n$ input strings represents a valid poker card, i.e.\ the string has
length $2$, its first character is one of \texttt{D C H S}, and its second
character is one of \texttt{A 2 3 4 5 6 7 8 9 T J Q K}.

\begin{center}
\begin{tabular}{ccc}
\toprule
Test case number & $n \leq$ & Special property \\
\midrule
$1$ & $1$ & A \\
$2 \sim 4$ & $52$ & A \\
$5 \sim 7$ & $52$ & B \\
$8 \sim 10$ & $52$ & none \\
\bottomrule
\end{tabular}
\end{center}

Special property A: it is guaranteed that the $n$ input cards are pairwise
distinct.

Special property B: it is guaranteed that all cards are given in increasing
order of rank, and cards of equal rank are given in the order diamonds,
clubs, hearts, spades.

\clearpage

\section*{Problem 2: Map Exploration (\begin{CJK}{UTF8}{gbsn}地图探险\end{CJK})}

Time limit: 1000 ms.
Memory limit: 512 MB.

Little A plans to go exploring in a jungle.
The geography of the jungle is very complicated, so to avoid getting lost he
first sends a robot to scout the way.

The map of the jungle can be represented by a table of characters with $n$
rows and $m$ columns.
We denote the coordinates of the position in row $i$ and column $j$ by
$(i, j)$ ($1 \leq i \leq n$, $1 \leq j \leq m$).
If the character at this position is \texttt{x}, then there is an obstacle
at this position and it cannot be passed.
Otherwise, if the character at this position is \texttt{.}, then this
position is empty ground and can be passed.

The state of the robot consists of two parts: a position and a direction.
The position is described by coordinates $(x, y)$
($1 \leq x \leq n$, $1 \leq y \leq m$), meaning that the robot is at the
position in row $x$ and column $y$ of the map.
The direction is given by an integer $d$ in $0 \sim 3$, where $d = 0$ means
facing east, $d = 1$ means facing south, $d = 2$ means facing west and
$d = 3$ means facing north.

Initially the robot is at position $(x_0, y_0)$ facing direction $d_0$.
\textbf{It is guaranteed that the initial position of the robot is empty
ground.}
The robot will then perform $k$ operations.
In each step the robot acts according to the following pattern:

\begin{enumerate}
\item Suppose the robot is currently at position $(x, y)$ facing
direction $d$.
Then the next position $(x', y')$ in its direction is defined as follows:
if $d = 0$, let $(x', y') = (x, y + 1)$;
if $d = 1$, let $(x', y') = (x + 1, y)$;
if $d = 2$, let $(x', y') = (x, y - 1)$;
if $d = 3$, let $(x', y') = (x - 1, y)$.
\item Next, the robot checks whether its next position lies inside the map
and is empty ground.
Concretely, it checks whether $(x', y')$ satisfies
$1 \leq x' \leq n$, $1 \leq y' \leq m$, and position $(x', y')$ is empty
ground.
If the condition holds, the robot moves one step forward.
Its new position becomes $(x', y')$ and its direction is unchanged.
If the condition does not hold, it performs a ``turn right'' operation.
That is, let $d' = (d + 1) \bmod 4$ (i.e.\ the remainder of $d + 1$
divided by $4$); the position stays the same, but the direction changes
from $d$ to $d'$.
\end{enumerate}

Little A wants to know how many positions on the map have been visited by
the robot (including the starting position) after the robot has performed
the $k$ operations.

\subsection*{Input}

\textbf{This problem has multiple test cases.}

The first line of input contains a positive integer $T$, the number of test
cases.

Then $T$ test cases follow, each in the following format:

The first line contains three positive integers $n, m, k$.
Here $n, m$ are the numbers of rows and columns of the map, and $k$ is the
number of operations the robot performs.

The second line contains two positive integers $x_0, y_0$ and a
non-negative integer $d_0$.

The next $n$ lines each contain a string of length $m$.
It is guaranteed that the strings contain only the two characters
\texttt{x} and \texttt{.}.
The $y$-th character of the $x$-th string corresponds to position $(x, y)$.
If this character is \texttt{x}, the position is an obstacle; otherwise it
is empty ground.
The data guarantee that the initial position of the robot is empty ground.

\subsection*{Output}

For each test case: output one line containing one positive integer, the
number of positions on the map visited by the robot (including the starting
position).

\subsection*{Sample 1 input}

\begin{verbatim}
2
1 5 4
1 1 2
....x
5 5 20
1 1 0
.....
.xxx.
.x.x.
..xx.
x....
\end{verbatim}

\subsection*{Sample 1 output}

\begin{verbatim}
3
13
\end{verbatim}

\subsection*{Explanation of sample 1}

This sample contains two test cases.
In the first test case the state of the robot changes as follows:
\begin{enumerate}
\item Initially the robot is at position $(1, 1)$ facing west (denoted by
the number $2$).
\item In the first step the robot finds that its next position $(1, 0)$ is
not inside the map, so it performs a ``turn right'' operation.
Its position is still $(1, 1)$, but it now faces north (denoted by the
number $3$).
\item In the second step the robot finds that its next position $(0, 1)$ is
not inside the map, so it again performs a ``turn right'' operation.
Its position is still $(1, 1)$, but it now faces east (denoted by the
number $0$).
\item In the third step the robot finds that its next position $(1, 2)$ is
inside the map and is empty ground.
Hence it moves one step east.
Its position becomes $(1, 2)$ and it still faces east.
\item In the fourth step the robot finds that its next position $(1, 3)$ is
inside the map and is empty ground.
Hence it moves one step east.
Its position becomes $(1, 3)$ and it still faces east.
\end{enumerate}

Therefore after four steps the robot has visited three positions, namely
$(1, 1), (1, 2), (1, 3)$.

In the second test case the operations performed by the robot in order are:
move east to $(1, 2)$, move east to $(1, 3)$, move east to $(1, 4)$,
move east to $(1, 5)$, turn right, move south to $(2, 5)$,
move south to $(3, 5)$, move south to $(4, 5)$, move south to $(5, 5)$,
turn right, move west to $(5, 4)$, move west to $(5, 3)$,
move west to $(5, 2)$, turn right, move north to $(4, 2)$, turn right,
turn right, move south to $(5, 2)$, turn right, turn right.

\subsection*{Sample 2}

See \texttt{explore/explore2.in} and \texttt{explore/explore2.ans} in the
contestant's directory.
(On Luogu these files are provided in the attachment \texttt{explore.zip}.)

This sample satisfies the constraints of test cases $3 \sim 4$.

\subsection*{Sample 3}

See \texttt{explore/explore3.in} and \texttt{explore/explore3.ans} in the
contestant's directory.

This sample satisfies the constraints of test case $5$.

\subsection*{Sample 4}

See \texttt{explore/explore4.in} and \texttt{explore/explore4.ans} in the
contestant's directory.

This sample satisfies the constraints of test case $6$.

\subsection*{Sample 5}

See \texttt{explore/explore5.in} and \texttt{explore/explore5.ans} in the
contestant's directory.

This sample satisfies the constraints of test cases $8 \sim 10$.

\subsection*{Constraints}

For all test data it is guaranteed that $1 \leq T \leq 5$,
$1 \leq n, m \leq 10^3$, $1 \leq k \leq 10^6$, $1 \leq x_0 \leq n$,
$1 \leq y_0 \leq m$, $0 \leq d_0 \leq 3$, and the starting position of the
robot is empty ground.

\begin{center}
\begin{tabular}{ccccc}
\toprule
Test case number & $n$ & $m$ & $k$ & Special property \\
\midrule
$1$ & $= 1$ & $\leq 2$ & $= 1$ & none \\
$2$ & $= 1$ & $\leq 2$ & $= 1$ & none \\
$3$ & $\leq 10^2$ & $\leq 10^2$ & $= 1$ & none \\
$4$ & $\leq 10^2$ & $\leq 10^2$ & $= 1$ & none \\
$5$ & $= 1$ & $\leq 10^3$ & $\leq 2 \times 10^3$ & all positions on the map are empty ground \\
$6$ & $= 1$ & $\leq 10^3$ & $\leq 2 \times 10^3$ & none \\
$7$ & $\leq 10^3$ & $\leq 10^3$ & $\leq 10^6$ & all positions on the map are empty ground \\
$8$ & $\leq 10^3$ & $\leq 10^3$ & $\leq 10^6$ & none \\
$9$ & $\leq 10^3$ & $\leq 10^3$ & $\leq 10^6$ & none \\
$10$ & $\leq 10^3$ & $\leq 10^3$ & $\leq 10^6$ & none \\
\bottomrule
\end{tabular}
\end{center}

\clearpage

\section*{Problem 3: Little Sticks (\begin{CJK}{UTF8}{gbsn}小木棍\end{CJK})}

Time limit: 1000 ms.
Memory limit: 512 MB.

Little S likes collecting little sticks.
After collecting $n$ little sticks of equal length, having nothing else to
do, he used them to form digits.
The way each digit is formed from little sticks is shown in the figure
below.

(The figure, omitted here, shows the digits $0$ to $9$ in the usual
seven-segment style; the numbers of sticks used by the digits
$0, 1, 2, 3, 4, 5, 6, 7, 8, 9$ are $6, 2, 5, 5, 4, 5, 6, 3, 7, 6$
respectively.)

Now Little S wants to form a \textbf{positive} integer satisfying the
following conditions:

\begin{itemize}
\item forming this number uses \textbf{exactly} $n$ little sticks;
\item the number formed has no leading $0$;
\item subject to the two conditions above, the number is as small as
possible.
\end{itemize}

Little S wants to know what this number is, but $n$ is very large and just
sorting out the sticks exhausted Little S, so you need to help him solve
this problem.
If no positive integer satisfies the conditions above, you need to output
$-1$ to report this.

\subsection*{Input}

This problem has multiple test cases.

The first line of input contains a positive integer $T$, the number of test
cases.

Then $T$ test cases follow, each in the following format:

One line containing an integer $n$, the number of sticks.

\subsection*{Output}

For each test case: output one line; if a positive integer satisfying the
requirements exists, output this number, otherwise output $-1$.

\subsection*{Sample 1 input}

\begin{verbatim}
5
1
2
3
6
18
\end{verbatim}

\subsection*{Sample 1 output}

\begin{verbatim}
-1
1
7
6
208
\end{verbatim}

\subsection*{Explanation of sample 1}

\begin{itemize}
\item For the first test case, no positive integer can be formed with
exactly one little stick, so output $-1$.
\item For the fourth test case, note that $0$ is not a valid answer.
Forming $9$, $41$ or $111$ each requires exactly $6$ little sticks, but
they are not the solutions with the smallest number formed.
\item For the fifth test case, forming $208$ requires $5 + 6 + 7 = 18$
little sticks.
It can be shown that forming any positive integer smaller than $208$
requires a number of little sticks different from $18$.
Note that although forming $006$ also requires $18$ little sticks, this
number has a leading zero and is therefore not a valid answer.
\end{itemize}

\subsection*{Constraints}

For all test data it is guaranteed that $1 \leq T \leq 50$,
$1 \leq n \leq 10^5$.

\begin{center}
\begin{tabular}{ccc}
\toprule
Test case number & $n \leq$ & Special property \\
\midrule
$1$ & $20$ & none \\
$2$ & $50$ & none \\
$3$ & $10^3$ & A \\
$4, 5$ & $10^5$ & A \\
$6$ & $10^3$ & B \\
$7, 8$ & $10^5$ & B \\
$9$ & $10^3$ & none \\
$10$ & $10^5$ & none \\
\bottomrule
\end{tabular}
\end{center}

Special property A: it is guaranteed that $n$ is a multiple of $7$ and
$n \geq 100$.

Special property B: it is guaranteed that there exists an integer $k$ with
$n = 7k + 1$, and $n \geq 100$.

\clearpage

\section*{Problem 4: Chain (\begin{CJK}{UTF8}{gbsn}接龙\end{CJK})}

Time limit: 2000 ms.
Memory limit: 512 MB.

Having grown used to playing idiom chains, Little J and his friends invented
a new chain rule.

In total $n$ people take part in this chain game; the $i$-th person
receives an integer sequence $S_i$ as his word bank.

A game consists of several rounds, and each round follows these rules:

\begin{itemize}
\item Some person $p$ among the $n$ people plays the round using his word
bank $S_p$.
If this is not the first round of the game, the person playing this round
must differ from the person of the previous round, but may coincide with the
person of the round before the previous one or any earlier round.
\item The player chooses a contiguous subsequence $A$ of $S_p$ of length in
$[2, k]$ as the \textbf{chain sequence} of this round, where $k$ is a given
constant.
If this is the first round of the game, $A$ must begin with the element $1$;
otherwise $A$ must begin with the last element of the chain sequence of the
previous round.
\begin{itemize}
\item A sequence $A$ is a contiguous subsequence of a sequence $S$ if and
only if $A$ can be obtained by deleting some elements (possibly none) from
the beginning and from the end of $S$.
\end{itemize}
\end{itemize}

To emphasise cooperation, Little J gives the $n$ participants $q$ tasks.
The $j$-th task requires the $n$ people to play one game in which exactly
$r_j$ rounds are played and the last element of the chain sequence of the
last round is exactly $c_j$.
To make sure that the tasks are feasible, Little J asks you to determine
whether each of the $q$ tasks can be completed, i.e.\ whether there exists
a possible course of the game satisfying the task's conditions.

\subsection*{Input}

\textbf{This problem has multiple test cases.}

The first line of input contains a positive integer $T$, the number of test
cases.

Then $T$ test cases follow, each in the following format:

The first line contains three integers $n, k, q$: the number of
participants, the upper bound on the length of a chain sequence, and the
number of tasks.

The next $n$ lines follow:

The $i$-th line contains $(l_i + 1)$ integers
$l_i, S_{i,1}, S_{i,2}, \dots, S_{i,l_i}$, where the first integer $l_i$ is
the length of the sequence $S_i$ and the next $l_i$ integers describe the
sequence $S_i$.

The next $q$ lines follow:

The $j$-th line contains two integers $r_j, c_j$ describing one task.

\subsection*{Output}

For each task: output one line containing one integer, 1 if the task can be
completed and 0 otherwise.

\newpage
\subsection*{Sample 1 input}

\begin{verbatim}
1
3 3 7
5 1 2 3 4 1
3 1 2 5
3 5 1 6
1 2
1 4
2 4
3 4
6 6
1 1
7 7
\end{verbatim}

\subsection*{Sample 1 output}

\begin{verbatim}
1
0
1
0
1
0
0
\end{verbatim}

\subsection*{Explanation of sample 1}

In what follows we write $\{A_i\} = \{A_1, A_2, \dots, A_r\}$ for all the
chain sequences of one game, and $\{p_i\} = \{p_1, p_2, \dots, p_r\}$ for
the numbers of the corresponding players.
Since all characters are single digits, for convenience we write sequences
directly as digit strings.

\begin{itemize}
\item For the first query, $p_1 = 1$, $A_1 = 12$ is a valid course of the
game satisfying the conditions.
\item For the second query, it can be shown that the task cannot be
completed.
Note that $p_1 = 1$, $A_1 = 1234$ is not a valid course of the game, since
then $|A_1| = 4 > k$.
\item For the third query, $\{p_i\} = \{2, 1\}$,
$\{A_i\} = \{12, 234\}$ is a valid course of the game satisfying the
conditions.
\item For the fourth query, it can be shown that the task cannot be
completed.
Note that $\{p_i\} = \{2, 1, 1\}$, $\{A_i\} = \{12, 23, 34\}$ is not a
valid course of the game: although all chain sequences have length at most
$k$, the second and third rounds are played by the same person, which is
not allowed.
\item For the fifth query, $\{p_i\} = \{1, 2, 3, 1, 2, 3\}$,
$\{A_i\} = \{12, 25, 51, 12, 25, 516\}$ is a valid course of the game
satisfying the conditions.
\item For the sixth query, it can be shown that the task cannot be
completed.
Note that every chain sequence must have length at least $2$, so
$A_1 = 1$ is not a valid course of the game.
\item For the seventh query, the character \texttt{7} does not occur in
anyone's word bank, so the task obviously cannot be completed.
\end{itemize}

\subsection*{Sample 2}

See \texttt{chain/chain2.in} and \texttt{chain/chain2.ans} in the
contestant's directory.
(On Luogu these files are provided in the attachment \texttt{chain.zip}.)

This sample satisfies the special property of test case 1.

\subsection*{Sample 3}

See \texttt{chain/chain3.in} and \texttt{chain/chain3.ans} in the
contestant's directory.

This sample satisfies the special property of test case 2.

\subsection*{Sample 4}

See \texttt{chain/chain4.in} and \texttt{chain/chain4.ans} in the
contestant's directory.

This sample satisfies special property A; moreover its first two test cases
satisfy $n \leq 1000$, $r \leq 10$, the sum of the lengths of all word
banks within one test case $\leq 2000$, and $q \leq 1000$.

\subsection*{Sample 5}

See \texttt{chain/chain5.in} and \texttt{chain/chain5.ans} in the
contestant's directory.

This sample satisfies special property B; moreover its first two test cases
satisfy $n \leq 1000$, $r \leq 10$, the sum of the lengths of all word
banks within one test case $\leq 2000$, and $q \leq 1000$.

\subsection*{Sample 6}

See \texttt{chain/chain6.in} and \texttt{chain/chain6.ans} in the
contestant's directory.

This sample satisfies special property C; moreover its first two test cases
satisfy $n \leq 1000$, $r \leq 10$, the sum of the lengths of all word
banks within one test case $\leq 2000$, and $q \leq 1000$.

\subsection*{Constraints}

For all test data it is guaranteed that:
\begin{itemize}
\item $1 \leq T \leq 5$;
\item $1 \leq n \leq 10^5$, $2 \leq k \leq 2 \times 10^5$,
$1 \leq q \leq 10^5$;
\item $1 \leq l_i \leq 2 \times 10^5$, $1 \leq S_{i,j} \leq 2 \times 10^5$;
\item $1 \leq r_j \leq 10^2$, $1 \leq c_j \leq 2 \times 10^5$;
\item letting $\sum l$ denote the sum of all $l_i$
\textbf{within one test case}, $\sum l \leq 2 \times 10^5$.
\end{itemize}

\begin{center}
\begin{tabular}{cccccc}
\toprule
Test case & $n \leq$ & $r \leq$ & $\sum l \leq$ & $q \leq$ & Special property \\
\midrule
$1$ & $10^3$ & $1$ & $2000$ & $10^3$ & none \\
$2, 3$ & $10$ & $5$ & $20$ & $10^2$ & none \\
$4, 5$ & $10^3$ & $10$ & $2000$ & $10^3$ & A \\
$6$ & $10^5$ & $10^2$ & $2 \times 10^5$ & $10^5$ & A \\
$7, 8$ & $10^3$ & $10$ & $2000$ & $10^3$ & B \\
$9, 10$ & $10^5$ & $10^2$ & $2 \times 10^5$ & $10^5$ & B \\
$11, 12$ & $10^3$ & $10$ & $2000$ & $10^3$ & C \\
$13, 14$ & $10^5$ & $10^2$ & $2 \times 10^5$ & $10^5$ & C \\
$15 \sim 17$ & $10^3$ & $10$ & $2000$ & $10^3$ & none \\
$18 \sim 20$ & $10^5$ & $10^2$ & $2 \times 10^5$ & $10^5$ & none \\
\bottomrule
\end{tabular}
\end{center}

Special property A: it is guaranteed that $k = 2 \times 10^5$.

Special property B: it is guaranteed that $k \leq 5$.

Special property C: it is guaranteed that within one test case, the total
number of occurrences of any one character in all word banks does not
exceed $5$.

\section*{Source}

The problems were translated from the official statements as hosted on
Luogu (the contest is not on QOJ):
\begin{itemize}
\item Problem 1: \texttt{https://www.luogu.com.cn/problem/P11227}
\item Problem 2: \texttt{https://www.luogu.com.cn/problem/P11228}
\item Problem 3: \texttt{https://www.luogu.com.cn/problem/P11229}
\item Problem 4: \texttt{https://www.luogu.com.cn/problem/P11230}
\end{itemize}
Time and memory limits are those shown on Luogu.
In the subtask tables, cells that Luogu displays as merged with the cell
above (``\verb|^|'') have been filled in with the repeated value.
The two figures of the original statements (the 52 cards of Problem 1 and
the stick digits of Problem 3) are replaced by textual descriptions.

\end{document}
